% ch2_answers.tex -- Quick check answers for Module 5, Chapter 2
\begin{chapteranswers}{2}{Circular Motion: Forces}
  % Q2.1 Ferris wheel
  \item One turn is the circumference, done in one period:
        $v = \dfrac{2\pi r}{T} = \dfrac{2\pi \times \qty{20}{m}}{\qty{80}{s}} = \qty{1.571}{m.s^{-1}}$.
        Angular velocity: $\omega = \dfrac{2\pi}{T} = \dfrac{2\pi}{80} = \qty{0.0785}{rad.s^{-1}}$.
        Answers: \final{$v = \qty{1.6}{m.s^{-1}}$} and \final{$\omega = \qty{0.079}{rad.s^{-1}}$}.
  % Q2.2 roundabout
  \item Yes. The speed stays the same, but the direction of motion keeps changing. Velocity
        is a vector, so the velocity changes, and a change in velocity is an acceleration.
        The acceleration points to the centre of the roundabout. (Friction from the road
        provides the force towards the centre.)
  % Q2.3 drum
  \item Radius: $r = \qty{48}{cm} \div 2 = \qty{24}{cm} = \qty{0.24}{m}$.
        Turns per second: $f = 1200 \div 60 = \qty{20}{Hz}$, so
        $T = \dfrac{1}{f} = \dfrac{1}{20} = \qty{0.050}{s}$.
        Speed: $v = \dfrac{2\pi r}{T} = \dfrac{2\pi \times 0.24}{0.050} = \qty{30.16}{m.s^{-1}}$.
        Acceleration: $a_\mathrm{c} = \dfrac{v^2}{r} = \dfrac{30.16^2}{0.24} = \qty{3.79e3}{m.s^{-2}}$,
        and $\dfrac{a_\mathrm{c}}{g} = \dfrac{3790}{9.8} = 387$.
        Answers: \final{$T = \qty{0.050}{s}$}, \final{$v = \qty{30}{m.s^{-1}}$},
        \final{$a_\mathrm{c} = \qty{3.8e3}{m.s^{-2}}$}, about 390 times $g$.
  % Q2.4 proportional
  \item $F_\mathrm{c} = \dfrac{mv^2}{r}$. (a) \qty{100}{km.h^{-1}} is twice the speed, so the
        force is $2^2 = 4$ times bigger: $4 \times 3.0 = \qty{12}{kN}$.
        (b) Half the radius at the same speed doubles the force: $2 \times 3.0 = \qty{6.0}{kN}$.
        Answers: \final{(a) \qty{12}{kN}} \ \final{(b) \qty{6.0}{kN}}.
  % Q2.5 real forces
  \item (a) Gravity: the Earth's pull on the Moon, towards the Earth's centre.
        (b) The normal force from the drum wall, pushing the clothes towards the axis.
        (c) Sideways (static) friction from the road on the tyres, towards the centre of
        the roundabout.
  % Q2.6 conical
  \item Angle: $\tan\theta = \dfrac{v^2}{rg} = \dfrac{1.4^2}{0.40 \times 9.8} = \dfrac{1.96}{3.92}
        = 0.5000$, so $\theta = \tan^{-1}(0.5000) = 26.57^\circ$ (keep the extra digits).
        Tension: divide both sides of the vertical equation $F_\mathrm{T}\cos\theta = mg$ by
        $\cos\theta$:
        $F_\mathrm{T} = \dfrac{0.50 \times 9.8}{\cos 26.57^\circ} = \dfrac{4.9}{0.8944} = \qty{5.48}{N}$.
        Check: $F_\mathrm{T}\sin\theta = 5.48 \times 0.4473 = \qty{2.45}{N}$ and
        $\dfrac{mv^2}{r} = \dfrac{0.50 \times 1.96}{0.40} = \qty{2.45}{N}$ \cmark.
        Answers: \final{$\theta = \ang{27}$} and \final{$F_\mathrm{T} = \qty{5.5}{N}$}.
  % Q2.7 velodrome
  \item $v = \sqrt{rg\tan\theta} = \sqrt{28 \times 9.8 \times \tan 43^\circ}
        = \sqrt{274.4 \times 0.9325} = \sqrt{255.9} = \qty{16.0}{m.s^{-1}}$.
        In \si{km.h^{-1}}: $16.0 \times 3.6 = 57.6$. Answer: \final{$\qty{58}{km.h^{-1}}$}.
  % Q2.8 swing
  \item Take $r$ = chain length = \qty{2.5}{m} (treat the child as a point at the end of
        the chains). At the lowest point the centre is above the child, so
        $F_\mathrm{T} - mg = \dfrac{mv^2}{r}$. Weight: $mg = 30 \times 9.8 = \qty{294}{N}$.
        Then $\dfrac{mv^2}{r} = \dfrac{30 \times 4.0^2}{2.5} = \qty{192}{N}$, so
        $F_\mathrm{T} = 192 + 294 = \qty{486}{N}$.
        Answer: \final{$F_\mathrm{T} = \qty{4.9e2}{N}$}. This is $\dfrac{486}{294} = 1.65$ times
        her weight, so she feels heavier.
  % Q2.9 flat bend (dry / ice)
  \item $v_\mathrm{max} = \sqrt{\mu g r}$.
        Dry: $\sqrt{0.75 \times 9.8 \times 65} = \sqrt{477.75} = \qty{21.86}{m.s^{-1}}$;
        $\times 3.6 = \qty{78.7}{km.h^{-1}}$.
        Ice: $\sqrt{0.10 \times 9.8 \times 65} = \sqrt{63.7} = \qty{7.98}{m.s^{-1}}$;
        $\times 3.6 = \qty{28.7}{km.h^{-1}}$.
        Answers: \final{dry \qty{79}{km.h^{-1}}} \ \final{ice \qty{29}{km.h^{-1}}}.
        The mass does not matter. The friction needed, $\dfrac{mv^2}{r}$, and the largest
        friction, $\mu mg$, both grow in the same way with $m$, so $m$ cancels.
  % Q2.10 rotor
  \item (a) Speed: $v = \dfrac{2\pi r}{T} = \dfrac{2\pi \times 2.0}{1.8} = \qty{6.9813}{m.s^{-1}}$.
        Riders stay up if $mg \le \mu\dfrac{mv^2}{r}$. Divide by $m$, multiply by $r$ and
        divide by $v^2$: $\mu \ge \dfrac{gr}{v^2}$. So
        $\mu_\mathrm{min} = \dfrac{gr}{v^2} = \dfrac{9.8 \times 2.0}{6.9813^2} = \dfrac{19.6}{48.74} = 0.402$.
        Answer: \final{$\mu = 0.40$} (no unit).
        (b) $F_\mathrm{N} = \dfrac{mv^2}{r}$, so the normal force gets smaller as the spin slows
        (half the speed gives a quarter of the force). The friction force stays the same:
        while the rider is stuck, the vertical forces balance, so $F_\mathrm{f} = mg$. Only the
        friction \emph{limit}, $\mu F_\mathrm{N}$, falls. When $\mu F_\mathrm{N}$ drops below $mg$,
        the rider slides down.
  % ---- 2.3 Key experiments (added by the coverage audit;
  % Q2.11 Part B: v^2 against r
  \item (a) It shows $v^2 \propto r$ when $F_\mathrm{c}$ and $m$ stay the same. The theory
        predicts this: rearranging $F_\mathrm{c} = \dfrac{mv^2}{r}$ gives
        $v^2 = \dfrac{F_\mathrm{c}}{m} \times r$. So at a fixed speed, a bigger radius needs a
        smaller force ($F_\mathrm{c} \propto \dfrac{1}{r}$).
        (b) The gradient is $\dfrac{F_\mathrm{c}}{m}$. Multiply both sides by $m$:
        $F_\mathrm{c} = \text{gradient} \times m = 57.5 \times 0.025 = \qty{1.44}{N}$.
        The hanging weight is $Mg = 0.150 \times 9.8 = \qty{1.47}{N}$.
        Percentage difference $= \dfrac{1.47 - 1.44}{1.47} \times 100\,\% = 2.2\,\%$ (from
        unrounded values). They agree within about 2\,\%, so the results support
        $F_\mathrm{c} = \dfrac{mv^2}{r}$.
        (c) From the line: $v^2 = 57.5 \times 0.80 = \qty{46.0}{m^2.s^{-2}}$, so
        $v = \sqrt{46.0} = \qty{6.78}{m.s^{-1}}$. Rearrange $v = \dfrac{2\pi r}{T}$:
        $T = \dfrac{2\pi r}{v} = \dfrac{2\pi \times 0.80}{6.78} = \qty{0.741}{s}$.
        Time for 20 turns $= 20 \times 0.741 = \qty{14.8}{s}$. Answer: \final{about \qty{15}{s}}.
  % Q2.12 Part C: v^2 against 1/m
  \item (a) Independent variable: the \textbf{stopper mass $m$} (the number of stoppers).
        Dependent variable: the \textbf{time for 20 turns} (it gives $T$, then $v$).
        Controlled variables (any two): the hanging mass $M$ (so $F_\mathrm{c}$ stays the same);
        the length $L$ (so $r$ stays the same); a horizontal circle; the number of turns timed.
        (b) $v^2$ against $\dfrac{1}{m}$. Rearranging $F_\mathrm{c} = \dfrac{mv^2}{r}$ gives
        $v^2 = F_\mathrm{c}\, r \times \dfrac{1}{m}$, so the gradient should be
        \final{$F_\mathrm{c}\, r$}.
        (c) Put $v = \dfrac{2\pi r}{T}$ into $F_\mathrm{c} = \dfrac{mv^2}{r}$:
        $F_\mathrm{c} = \dfrac{m (2\pi r)^2}{T^2\, r} = \dfrac{4\pi^2 m r}{T^2}$. Multiply both
        sides by $T^2$ and divide both sides by $F_\mathrm{c}$:
        $T^2 = \dfrac{4\pi^2 m r}{F_\mathrm{c}}$. Here $F_\mathrm{c}$ and $r$ are fixed, so
        $T^2 \propto m$ and $T \propto \sqrt{m}$. Doubling $m$ multiplies $T$ by
        $\sqrt{2} = 1.414$: $T = 0.60 \times 1.414 = \qty{0.849}{s}$.
        Answer: \final{$T = \qty{0.85}{s}$}. The heavier stopper goes round more slowly.
  % Q2.13 Part A: F_c against v or v^2
  \item From the table: at $F_\mathrm{c} = \qty{0.98}{N}$, $v^2 = \qty{36.0}{m^2.s^{-2}}$, so
        $v = \sqrt{36.0} = \qty{6.00}{m.s^{-1}}$. At $F_\mathrm{c} = \qty{1.96}{N}$,
        $v^2 = \qty{72.1}{m^2.s^{-2}}$, so $v = \sqrt{72.1} = \qty{8.49}{m.s^{-1}}$.
        Doubling $F_\mathrm{c}$ multiplies $v$ by $\dfrac{8.49}{6.00} = 1.42$, which is about
        $\sqrt{2}$. Answer: \final{$v$ is multiplied by about 1.4}, not by 2.
        If $F_\mathrm{c}$ were proportional to $v$, the speed would double to
        \qty{12.0}{m.s^{-1}}. It does not, so $F_\mathrm{c}$ against $v$ is a curve.
        But $\dfrac{F_\mathrm{c}}{v^2}$ stays the same: $\dfrac{0.98}{36.0} = 0.0272$ and
        $\dfrac{1.96}{72.1} = 0.0272$ (unit \si{kg.m^{-1}}). So $F_\mathrm{c} \propto v^2$, and
        $F_\mathrm{c}$ against $v^2$ is a straight line through the origin. Its gradient is
        $\dfrac{m}{r}$.
\end{chapteranswers}
